% KKchemstruct user manual (English)
%   lualatex KKchemstruct-manual-en.tex   (twice)
% Run twice: \KKchemmark and \KKhbond need a second pass to settle.
\documentclass[paper=a4paper,fontsize=10pt]{jlreq}
\usepackage{KKchemstruct}

\title{KKchemstruct User Manual}
\author{KKTeX}
\date{2 August 2026\quad v1.0.0}

\begin{document}
\maketitle
\tableofcontents

\section{Overview}

\textsf{KKchemstruct} is a set of macros for typesetting the structural formulas
of high-school chemistry (organic and inorganic). The drawing itself is done by
\textsf{chemfig}; this package configures its dimensions, line widths and
orientations to match the look of printed textbook figures, and adds a short
notation for substituted benzenes, condensed formulas, reaction schemes and
annotations.

Every dimension is defined in \verb|em|, so the figures follow the body font
size. There is no need to rewrite the numbers for each paper size.

\subsection{Requirements}

\begin{itemize}
\item Engine: LuaLaTeX, or pLaTeX + dvipdfmx
\item Packages: \textsf{chemfig}, \textsf{etoolbox}, \textsf{xkeyval}
      (all part of TeX Live)
\end{itemize}

\subsection{Loading}

Put \verb|KKchemstruct.sty| on the TeX search path or next to your document and
write

\begin{verbatim}
\usepackage{KKchemstruct}
\end{verbatim}

There are no package options.

\subsection{When two compilations are needed}

\verb|\KKchemmark| and \verb|\KKhbond| refer to TikZ nodes, so they need two
compilations to settle. On the first pass the frame or the dotted line appears
in the wrong place; the second pass puts it right.

\section{Benzene rings}

\subsection{\texttt{\char`\\KKbenzene}}

Substituted benzenes are written with \verb|\KKbenzene|. Positions are numbered
1 to 6 clockwise, starting from the top vertex.

\begin{verbatim}
\KKbenzene[2=OH,3=COOH]
\end{verbatim}

\begin{center}
\KKbenzene[2=OH,3=COOH]\qquad
\KKbenzene[1=COOH,4=COOH]\qquad
\KKbenzene[]
\end{center}

Substituents are set in math mode, so subscripts such as \verb|COOCH_3| work as
written. Positions you leave out produce no branch at all.

\subsection{Options}

\begin{description}
\item[\texttt{1}--\texttt{6}] the substituent at each position.
\item[\texttt{angle}] rotation of the whole ring, in degrees. The default is no
      rotation (vertices at top and bottom).
\item[\texttt{sep}] the side length (distance between atom centres) of this ring only.
\item[\texttt{kekule=b}] flips the inner double-bond strokes to the upper-right,
      bottom and upper-left sides. Use it when a substituent collides with a stroke.
\end{description}

\begin{verbatim}
\KKbenzene[1=OH,kekule=b]      \KKbenzene[2=OH,angle=30,sep=1.0em]
\end{verbatim}

\begin{center}
\KKbenzene[1=OH,kekule=b]\qquad
\KKbenzene[2=OH,angle=30,sep=1.0em]
\end{center}

\subsection{Rings set inline}

To put a ring in running text, use the flat-top form, whose top and bottom sides
are horizontal. It bonds at the left and right vertices, so it sits naturally in
the middle of a chain.

\begin{description}
\item[\texttt{\char`\\KKphenylring}] the bare ring.
\item[\texttt{\char`\\KKphenyl\{...\}}] a ring with a substituent to its right.
\item[\texttt{\char`\\KKphenylene\{...\}}] adds a bond stub on the left as well,
      so the ring can be inserted into a chain.
\end{description}

A ring such as \KKphenylring{} in running text does not disturb the leading.
A substituent can be attached on the right, as in \KKphenyl{COOH}, and a ring can
be dropped into a chain: \ck{CH_3(CH_2)_{11}}\KKphenylene{SO_3H}.

The amount by which such rings are raised above the baseline is set with the
\verb|inline raise| key:

\begin{verbatim}
\KKchemsetup{inline raise=0.35ex}
\end{verbatim}

\subsection{Fused rings (\texttt{\char`\\KKchemring})}

For rings other than benzene, and for fused rings, pass raw \textsf{chemfig}
notation to \verb|\KKchemring|. Only the ring dimensions (constant distance
between atom centres) are applied. The optional first argument takes
\textsf{chemfig} keys.

\begin{verbatim}
\KKchemring[atom sep=1.75em]{*6(=-*5(-C(=O)-O-C(=O)-)=-=-)}
\end{verbatim}

\begin{center}
\KKphthalicanhydride\qquad\KKmaleicanhydride
\end{center}

Rings that carry C or O labels on their vertices need a longer side than
skeleton-only rings, or the labels crowd each other. The default for those is
\verb|label ring sep| (1.75em).

\section{Condensed formulas}

\subsection{\texttt{\char`\\KKchemchain}}

Formulas in which labels run horizontally, such as
\verb|CH_2-O-C(=[2]O)-R_1|, are written with \verb|\KKchemchain|. Here the
\emph{visible} bond length is kept constant, so the line is not squeezed after a
wide label such as \verb|CH_2|.

\begin{verbatim}
\KKchemchain{CH_2(-OH)-[6,1.15]CH(-OH)-[6,1.15]CH_2-OH}
\end{verbatim}

\begin{center}
\KKglycerol\qquad\KKtriglyceride{R_1}{R_2}{R_3}
\end{center}

Angles use \textsf{chemfig} numbering in steps of 45 degrees (0 = right, 2 = up,
4 = left, 6 = down). The number after the angle is a length coefficient.

\subsection{Common fragments}

Ready-made pieces for use inside a chain. They are \textsf{chemfig} code
fragments, so they belong in the argument of \verb|\KKchemchain|.

\begin{description}
\item[\texttt{\char`\\KKcarbonyl}] an upward C=O.
\item[\texttt{\char`\\KKcarbonyldown}] a downward C=O.
\item[\texttt{\char`\\KKdative\{<angle>\}}] a dative bond arrow (S$\to$O, P$\to$O, N$\to$O).
\end{description}

Vertical bonds are drawn longer than horizontal ones: without allowing for the
height of the labels and for the arrowhead, they look cramped.

\begin{verbatim}
\KKchemchain{H-O-S(\KKdative{2}O)(\KKdative{6}O)-O-H}
\end{verbatim}

\begin{center}
\KKsulfuricacid\qquad\KKnitricacid\qquad\KKphosphoricacid
\end{center}

\section{Reaction schemes}

\subsection{\texttt{\char`\\KKscheme}}

A scheme is written by putting \textsf{chemfig} \verb|\arrow| commands into the
argument of \verb|\KKscheme|. The arrow syntax is \textsf{chemfig}'s own.

\begin{verbatim}
\KKscheme{%
  \KKname{\KKsalicylicacid}{salicylic acid}
  \arrow{->[\ck{CH_3OH}][esterification]}
  \KKname{\KKmethylsalicylate}{methyl salicylate}}
\end{verbatim}

\begingroup
\KKchemsetup{compound sep=9em}
\KKscheme{%
  \KKname{\KKsalicylicacid}{salicylic acid}
  \arrow{->[\ck{CH_3OH}][esterification]}
  \KKname{\KKmethylsalicylate}{methyl salicylate}}
\endgroup

The optional first argument takes \verb|\schemestart| keys
(\verb|arrow angle|, \verb|arrow coeff|, \verb|arrow style|).

The default distance between compounds, \verb|compound sep|, is 5.5em, which
suits short labels such as chemical formulas. A long word set over the arrow --
\textsf{esterification} above -- is wider than that, and would run into the
compound on its left; the example is therefore wrapped in

\begin{verbatim}
\begingroup\KKchemsetup{compound sep=9em} ... \endgroup
\end{verbatim}

\subsection{\texttt{\char`\\KKbranchscheme}}

A figure in which one substrate branches into two products is written with
\verb|\KKbranchscheme|, which takes seven arguments. Each arrow rises at an
angle and then runs horizontally; the labels sit above and below the horizontal
part.

\begin{verbatim}
\KKbranchscheme
  {<substrate>}
  {<upper reagent>}{<upper reaction>}{<upper product>}
  {<lower reagent>}{<lower reaction>}{<lower product>}
\end{verbatim}

\KKbranchscheme
  {\KKname{\KKsalicylicacid}{salicylic acid}}
  {\ck{CH_3OH}}{esterification}{\KKname{\KKmethylsalicylate}{methyl salicylate}}
  {\ck{(CH_3CO)_2O}}{acetylation}{\KKname{\KKaspirin}{acetylsalicylic acid}}

The angle and the length of the branches are set with \verb|branch angle|
(default 20 degrees) and \verb|branch coeff| (default 2.7).

\subsection{Compound names (\texttt{\char`\\KKname})}

\verb|\KKname{<structure>}{<name>}| centres a name under a structure. The font is
set with \verb|name font|.

\section{Annotations}

\subsection{Marking a leaving group (\texttt{\char`\\KKchemmark})}

The \verb|-OH| and \verb|-H| lost in a condensation can be enclosed in a dotted
rounded frame. Tag the atoms with \textsf{chemfig}'s \verb|@{name}| and give the
two tags as opposite corners.

\begin{verbatim}
\KKchemchain{R-C\KKcarbonyl-@{s}OH}\ +\ \KKchemchain{@{e}H-O-R'}%
\KKchemmark{s}{e}
\end{verbatim}

\begin{center}
\KKchemchain{R-C\KKcarbonyl-@{s}OH}\ +\ \KKchemchain{@{e}H-O-R'}%
\KKchemmark{s}{e}
\end{center}

The optional first argument takes TikZ keys, for example to change the colour.

\subsection{Hydrogen bonds (\texttt{\char`\\KKhbond})}

Draws the dotted line of an intramolecular hydrogen bond between two tags.

\begin{verbatim}
\KKbenzene[2=@{oh}OH,3=C@{co}OOH]\KKhbond{oh}{co}
\end{verbatim}

\begin{center}
\KKbenzene[2=@{oh}OH,3=C@{co}OOH]\KKhbond{oh}{co}
\end{center}

\subsection{Crossing out (\texttt{\char`\\KKchemcross})}

Superimposes a heavy cross to show that a reaction does not take place. The
optional first argument sets the colour (default \verb|black!35|).

\begin{center}
\KKchemcross{\KKphenol}
\end{center}

\section{Chemical formula macros}

\textsf{chemfig} sets atom labels in math mode, so subscripted formulas such as
\verb|CH_3| work as written inside a structure. Arrow labels, however, are not in
math mode -- that is how \textsf{chemfig} works. Wrap those in one of the
following.

\begin{description}
\item[\texttt{\char`\\KKchemformula\{...\}}] sets a formula with \verb|\mathrm|.
\item[\texttt{\char`\\ck\{...\}}] a short alias for the above. It exists only to save
      keystrokes, so use \verb|\KKchemformula| in documents where a two-letter
      name would clash.
\item[\texttt{\char`\\ckm} / \texttt{\char`\\ckp}] the superscript $-$ and $+$ of an ion.
\end{description}

Use them as in \ck{CH_3COOH}, \ck{RCOO\ckm} and \ck{Na\ckp}. Labels in prose need
no wrapper.

\section{Changing the dimensions}

\subsection{\texttt{\char`\\KKchemsetup}}

All dimensions are changed through \verb|\KKchemsetup|.

\begin{verbatim}
\KKchemsetup{ring sep=2.0em,bond width=0.09em,double bond sep=0.4em}
\end{verbatim}

\begingroup
\KKchemsetup{ring sep=2.0em,bond width=0.09em,double bond sep=0.4em}
\begin{center}
\KKsalicylicacid\qquad\KKphenyl{COOH}
\end{center}
\endgroup

Call it inside a group and the change stays in that group. The example above is
wrapped in \verb|\begingroup| ... \verb|\endgroup|, so here the original
dimensions are back in force: \KKsalicylicacid

\subsection{List of keys}

\begin{description}
\item[\texttt{ring sep} (1.35em)] side of a skeleton-only ring (distance between atom centres).
\item[\texttt{label ring sep} (1.75em)] side of a ring whose vertices carry C or O labels.
\item[\texttt{chain sep} (0.75em)] bond length in condensed formulas; the visible
      line is shorter than this by the gap at each end.
\item[\texttt{vertical coeff} (1.15)] how much vertical bonds are lengthened.
\item[\texttt{dative coeff} (1.5)] how much dative-bond arrows are lengthened.
\item[\texttt{mid bond coeff} (1.28)] coefficient of the middle horizontal bond in a triglyceride.
\item[\texttt{bond offset} (0.13em)] gap between a label and a bond line.
\item[\texttt{double bond sep} (0.26em)] distance between the two lines of a double bond.
\item[\texttt{bond width} (0.05em)] thickness of bond lines.
\item[\texttt{arrow label sep} (0.30em)] gap between a reaction arrow and its label.
\item[\texttt{arrow offset} (0.8em)] gap between a compound and the end of an arrow.
\item[\texttt{compound sep} (5.5em)] distance between compounds in a scheme.
\item[\texttt{kink} (1.8em)] the rise of a branching arrow.
\item[\texttt{branch angle} (20) / \texttt{branch coeff} (2.7)] angle and length of branching arrows.
\item[\texttt{inline raise} (0.35ex)] how far an inline ring is raised.
\item[\texttt{name font} (\texttt{\char`\\normalsize})] font of the name set by \verb|\KKname|.
\end{description}

\section{Predefined compounds}

Frequently used structures can be called directly. Each is just a combination of
the macros described above, so you can add your own by following the pattern at
the end of the \verb|.sty| file.

\begin{description}
\item[Aromatics] \verb|\KKsalicylicacid|, \verb|\KKmethylsalicylate|, \verb|\KKaspirin|,
      \verb|\KKbenzoicacid|, \verb|\KKphenol|, \verb|\KKphthalicacid|,
      \verb|\KKisophthalicacid|, \verb|\KKterephthalicacid|
\item[Anhydrides and geometric isomers] \verb|\KKphthalicanhydride|,
      \verb|\KKmaleicanhydride|,\\
      \verb|\KKmaleicacid|, \verb|\KKfumaricacid|
\item[Oxoacids] \verb|\KKnitricacid|, \verb|\KKsulfuricacid|, \verb|\KKphosphoricacid|
\item[Fats] \verb|\KKglycerol|, \verb|\KKtriglyceride{R_1}{R_2}{R_3}|, \verb|\KKestergeneric|
\end{description}

\begin{center}
\KKmaleicacid\qquad\KKfumaricacid
\end{center}

\section{Pitfalls}

\begin{description}
\item[Parentheses make branches] In \textsf{chemfig}, \verb|(| and \verb|)| are the
      branch notation. Passing something like \verb|CH_3(CH_2)_{11}| straight to
      \verb|\KKchemchain| makes it a branch. Write formulas that are not drawn as
      structures with \verb|\ck{...}| instead.
\item[Arrow labels are not in math mode] As above, wrap them in \verb|\ck{...}|.
\item[Annotations need two passes] \verb|\KKchemmark| and \verb|\KKhbond| are not
      positioned on the first compilation.
\item[Rings and chains fix different lengths] Rings keep the distance between atom
      centres constant, which is what makes a regular hexagon; chains keep the
      visible line length constant, so wide labels do not squeeze the line. Use
      \verb|\KKbenzene| / \verb|\KKchemring| for rings and \verb|\KKchemchain| for
      chains. Mixing them is safe: each macro switches the dimensions itself.
\end{description}

\section{License}

MIT License. See \verb|LICENSE.md| in the repository for the full text.

\end{document}
